% Source: sources/returns/S6_04_MAPS_MANIFOLDS/sections/manifolds.tex:93–127
\paragraph{Theorem D.1.3 (Proved: the rational homology of the cover).}
For every admissible $p$, the rational homology of $\widetilde X_p$ has
ranks
\[
 (1,0,d-1,0,d-1,0,1).
\]
The deck group on either $H^2(\widetilde X_p;\mathbb C)$ or
$H^4(\widetilde X_p;\mathbb C)$ is the direct sum of its nontrivial
characters, each with multiplicity one. This statement does not
assert the absence of integral torsion.
\begin{proof}
Let $\lambda\ne1$ be a $d$th root of unity and let
$\mathcal L_\lambda$ be the associated rank-one complex local system
on $X_p$. Its restriction to a regular fibre is nontrivial on
$\widehat\gamma$. The cellular cochain complex of a torus with this
local system is a Koszul complex containing the invertible scalar
$\lambda-1$ and is therefore contractible. On a reduced multiple
fibre the finite torus cover has the same property. Transfer, or
averaging over the finite deck group in characteristic zero, shows
that its cohomology also vanishes. The regular region and all three
boundary overlaps consequently have zero cohomology with these
coefficients, as follows either by their bundle spectral sequences
or by contraction of the fibre cochain complex. Mayer--Vietoris
therefore identifies $H^*(X_p;\mathcal L_\lambda)$ with
$H^*(W;\mathcal L_\lambda|_W)$.
In the preceding normalization calculation substitute $t=\lambda$.
The degree-zero row is acyclic, the matrix $B(\lambda)$ has rank three
and one-dimensional kernel, and the degree-four row is one-dimensional.
Thus this local system contributes one copy of $\mathbb C$ in degrees
two and four and nothing else. The trivial character contributes
$H^*(X_p;\mathbb C)$, which is concentrated in degrees zero and six by
\textup{[AC26, Theorem 7.22]}. The decomposition of the regular
representation of $\mathbb Z/d$ into its characters computes the
cohomology of the covering and proves the claim.
\end{proof}
